


Optimal Control (OC) provides a principled framework for designing a control law by specifiying it to be the minimizer of a specified performance index. 
However, the performance of such controller in the contextual environment critically depends on the choice of the cost function, which in general trade-offs between the control effort and the feedback system's performance. In practice, accurate specification of such cost function is often nontrivial; it requires expertise and may not reflect the true objectives of the expert's decisions.
Inverse optimal control (IOC), also known as Inverse Reinforcement Learning (IRL), addresses this issue by inferring the underlying cost function of an expert from demonstration data. 
Therefore, it not only provides valuable insights for controller design, but also further enables the reconstruction of the corresponding closed-loop policy \citep{dey2023inverse, mombaur2010human}.


Compared with Imitation Learning (IL), which directly imitates the observed expert behavior \citep{zare2024survey}, IOC assumes that the expert's cost function follows a known structure, such as a linear combination of predefined feature functions \citep{ziebart2008maximum,zhang2019inverse,garrabe2025convex}.
This assumption restricts the optimal policy's possible structure and provides stronger robustness and generalizability.
Motivated by this, IOC methods have been successfully applied to modeling the decision criteria in various domains such as autonomous driving \citep{likmeta2021dealing}, human sensorimotor \citep{Schultheis2021Inverse}, and animal behaviours \citep{Ashwood2022Dynamic}.
However, for nonlinear stochastic systems, the corresponding IOC problems are challenging due to the absence of explicit expressions for the value functions.
This leads to ``coupled-grey-black-box" identification problems, compared to the grey-box case in Linear Quadratic (LQ) settings \citep{zhang2019inverse,zhang2019inverseCDC,li2018convex,zhang2021inverse,zhang2022statistically,yu2021system, Menner2018Convex, rickenbach2024inverse, guo2023imitation, zhang2024statistically, zhang2024inverse}.
Such problem still remains even when the cost function structure is known a priori. 
Moreover, the trajectory observations are usually contaminated by noise in practice. In particular, the noise could be either the observation noise introduced by the sensor measurements, or the process noise that is intrinsically modelled by the Markov transition kernel of the nonlinear stochastic system.
Therefore,  a key property that an IOC estimator should have to guarantee the robustness is consistency.
Nevertheless, if the IOC estimator is an ``M-estimator" \citep[Chap.~5]{van1998asymptotic}, consistency is typically only guaranteed at the globally optimal solution of the underlying optimization problem. Therefore,  it is also desirable that the this optimization is convex, as to avoid issues with local minima and thus achieve consistency in practice.


\subsection{Related works}
One way to avoid the aforementioned ``coupled-grey-black-box" issue is to identify a cost function that satisfies the first-order necessary optimality conditions given by Pontryagin's Maximum Principle (PMP) \citep{Liang2023DataDriven}.
However, such algorithms only recovers a cost function to which the observed trajectories is only locally optimal, and hence we might not be able to guarantee the feedback system's performance when using the control policy recovered from the estimated cost function.
Another straightforward method is to minimize the discrepancy between the expert's policy and the optimal policy induced by the candidate cost functions \citep{abbeel2004apprenticeship, ziebart2008maximum, Ho2016GAIL}.
Such methods often require repeatedly solving the entire or part of the forward optimal control problem, which makes them computationally demanding and less scalable to complex systems.
Nevertheless, by exploiting the specific structures of the system dynamics or the cost function, several studies have avoided the issue of solving the forward problem repeatedly.
In particular, bilinear \citep{fernandez2025estimating}, affine \citep{lian2022inverse, Jean2019Injectivity, Rodrigues2022Inverse}, and polynomial dynamical systems \citep{pauwels2016linear} have been investigated, where certain problem structures enable tractable reformulations. Yet more general nonlinear settings are not considered therein.
Furthermore, another line of research incorporates statistical learning principles.
Inspired by maximum-entropy regularization, a family of methods recovers the cost function by maximum-likelihood estimation under stochastic expert policies \citep{ziebart2008maximum, dvijotham2010inverse, garrabe2025convex, Mehr2023Maximum, Guo2021Deep}.
In addition, \citep{ziebart2008maximum, zhou2018Infinite} discretize the state and control spaces, while others represent the stochastic policies with exponential-family distributions \citep{dvijotham2010inverse, garrabe2025convex, Mehr2023Maximum} or deep neural networks \citep{Guo2021Deep}.
Nevertheless, the impact of these approximations on estimation accuracy remains theoretically unquantified.

\subsection{The contributions}
To address the above challenges, we propose an IOC algorithm for discrete-time infinite-horizon nonlinear stochastic systems with a discounted-cost, where the running cost is assumed to be a linear combination of known continuous feature functions.
In particular, the method recovers the underlying cost function by solving a Sum-Of-Square (SOS) programming problem, which is convex and well-studied. 
Compared to existing approaches, our method
(i) handles the nonlinear and stochastic dynamics without restrictive structural assumptions;
(ii) avoids solving the forward optimal control problem repeatedly; and 
(iii) provides theoretical guarantees of asymptotic and statistical consistency, which have not been established simultaneously in existing works.
More precisely, our key contributions are as follows:
\begin{enumerate}[label=\roman*.]
    \item We reformulate the nonlinear, infinite-horizon IOC problem as an infinite dimensional linear programming problem, which minimizes the violation of optimality conditions for the observed trajectories. We further show that samples of finite-time observation suffice to identify the cost function parameters. By applying a polynomial approximation, we obtain a finite dimensional SOS program that can be solved numerically. Such method also exempts us from local minima issues.
    \item We establish asymptotic and statistical consistency of the proposed estimator. Namely, as the polynomial approximation order and the number of demonstration trajectories increases, the optimality gap between the expert policy and the estimated optimal policy vanishes, and the estimated parameters approach the true ones whenever the solution is unique. Moreover, the underlying optimal control policy can be further consistently recovered, e.g., by integrating the estimated optimality condition into ``behaviour cloning''.
    \item Numerical experiments are conducted to validate the theoretical consistency of the proposed method, demonstrate its estimation accuracy, and evaluate its robustness and generalizability.
\end{enumerate}


The remainder of this paper is organized as follows.
Sec.~\ref{sec:problem_formulation} introduces the problem setting, states several mild assumptions on the considered system, and formulates the IOC problem for discrete-time nonlinear stochastic systems.
Next, we present the proposed approach in Sec.~\ref{sec:alg_construct}. In particular, the unknown functions in the IOC formulation are approximated by polynomials, and the problem is converted into a finite-dimensional SOS program --- this is the finite-dimensional estimator.
Moreover, the theoretical guarantees, including asymptotic and statistical consistency of the estimator, are established in Sec.~\ref{sec:consistency}.
In addition, we briefly describes the SOS optimization procedure used in practice in Sec.~\ref{sec:practical_implementation}. 
In Sec.~\ref{sec:numerical-expe}, we evaluate perform  numerical experiments and the estimator on three benchmark systems: a LQ regulator, a temperature-control process, and an inverted pendulum. In the inverted pendulum problem, we also compare the controllers obtained via IOC with ``behaviour cloning'' (see, e.g., \citep{sammut2017behavioral, codevilla2019exploring, torabi2018behavioralcloningobservation}) and show that IOC gives controllers that are more robust to distributional shifts and errors in the modelled dynamics.
Finally, the paper is concluded in Sec.~\ref{sec:conclusions} with some discussion and outlooks.

\textit{Notation}: 
$\mR_+$ denotes the non-negative real numbers. $C(n,k)$ denotes the binomial coefficient (``$n$ choose $k$").
We denote $\mathcal{M}(A)$ as the space of all finite signed measures on $A$, and let $\support(\mu)$ denote the support of $\mu\in\mathcal{M}(A)$, while $\mathscr{F}(A)$ denotes the space of all bounded measurable functions on $A$. In addition, let $\contfun(A)$ denote the set of continuous functions whose domain is $A$. 
Furthermore, we denote $\|f\|_{p,A}:=\left(\int_A\left|f(x)\right|^p\mathrm{d}x\right)^{\frac{1}{p}}$ as the $p$-norm of the bounded function $f$ on the set $A\subset \mR^n$. In both this and the previous notation, we sometimes drop the set $A$ when there is no risk of confusion.
For vector valued function $f:\mR^n\rightarrow\mR^n$, $\|f\|_{\infty,A}:=\max_{i=1:n}\|f_i\|_{\infty,A}$, where $f_i$ is the $i$-th element of $f$.
Similarly, $\|v\|_p:=\left(\sum_{i=1}^n|v_i|^p\right)^{\frac{1}{p}}$ for $v=[v_1,\ldots,v_n]^T\in\mR^n$.
On the other hand, we denote $\|\mu\|_{\tv}:=\sup_{\|v\|_\infty\leq 1}\int v(x) \mu(\mathrm{d}x)$ as the total variation norm.
We also denote $\lip(\beta)\subset \mathcal{C}(A)$ as the set of all Lipschitz continuous functions on $A$ that has the Lipschitz constant $\beta$.
We further denote $\mR[x]$ as the ring of multivariate polynomials of $x\in\mR^n$ and $\Sigma^2$ as the SOS polynomials. More specifically, $p\in\Sigma^2$ if there exists an integer $m$ and $u_i\in\mR[x]$, such that $p=\sum_{i=1}^m u_i^2$.
Moreover, we define the Cartesian product as $\times_{i=1}^n A_i := A_1\times \ldots \times A_n$.
$\langle\cdot,\cdot\rangle$ denotes the dual bracket between the dual pair $(\cdot,\cdot)$.
Moreover, for linear operator $\mathscr{G}$, we let $\mathscr{G}^*$ denote its adjoint.
For any class $B\subset \mR^n$, the diameter is defined as $\diam(B):=\sup_{x,y\in B}\|x-y\|_2$.
In addition, $\indicator_{\Gamma}(x)$ denotes the indicator function of $x$ on set $\Gamma$, namely,
\[
\indicator_\Gamma(x) = 
\begin{cases}
1 &\mbox{if } x\in\Gamma,\\
0 &\mbox{otherwise}.
\end{cases}
\]
Moreover, we use $\textbf{\textit{italic bold font}}$ to denote stochastic elements, and $\mfx \sim \mu$ means that $\mfx$ has distribution $\mu$. Specifically, $\mathcal{N}(b, \sigma^2)$ denotes normal distribution with mean $b$ and variance $\sigma^2$ and $\mathcal{N}_C(b, \sigma^2)$ denotes the truncated normal distribution on the hyper-cube $C$. $\mathcal{U}(A)$ denotes uniform distribution on $A$.
